\documentclass[11pt,a4paper]{article}
\usepackage[italian]{babel}
\usepackage[margin=2.5cm]{geometry}
\usepackage{parskip}
\usepackage{amsmath, amssymb, amsthm}
\usepackage[most]{tcolorbox}
\usepackage{array}
\usepackage{tikz}
\usetikzlibrary{decorations.pathreplacing}
\usepackage{hyperref}

% --- Riquadri con bordo nero, senza numero ---
% Uso: \begin{definizione} ... \end{definizione}
%      \begin{teorema}[Nome] ... \end{teorema}  (il nome è facoltativo)
\tcbset{nota/.style={colback=white, colframe=black, boxrule=0.5pt,
  sharp corners}}
\NewTColorBox{definizione}{}{nota,
  before upper={\textbf{Definizione.}\ }}
\NewTColorBox{teorema}{o}{nota,
  before upper={\textbf{Teorema\IfValueT{#1}{ (#1)}.}\ }}
\NewTColorBox{proposizione}{o}{nota,
  before upper={\textbf{Proposizione\IfValueT{#1}{ (#1)}.}\ }}
\NewTColorBox{assioma}{o}{nota, boxrule=1pt,
  before upper={\textbf{Assioma\IfValueT{#1}{ (#1)}.}\ }}

% --- Ambienti semplici (senza riquadro) ---
\theoremstyle{definition}
\newtheorem*{esempio}{Esempio}
\newtheorem*{osservazione}{Osservazione}
\newtheorem*{esercizio}{Esercizio}

% V e F nelle tavole di verità
\newcommand{\V}{\textsf{V}}
\newcommand{\F}{\textsf{F}}

\title{Funzioni in due variabili}
\author{Marco Cerbella}
\date{Lezione 4 -- 5 ottobre 2026}

\begin{document}
\maketitle

\section{Funzioni in due variabili}

\begin{definizione}
Si dice \emph{funzione di due variabili} una applicazione che ha come dominio un sottoinsieme di $\mathbb{R}^2$ e come codominio un sottoinsieme di $\mathbb{R}$:
\[
f : D \subseteq \mathbb{R}^2 \to \mathbb{R}, \qquad (x, y) \mapsto f(x, y).
\]
\end{definizione}

\begin{itemize}
    \item \emph{Grafico}: $\mathrm{Gr}(f) = \{(x, y, z) \in \mathbb{R}^2 \times \mathbb{R} : z = f(x, y),\ (x, y) \in D\}$ (una superficie nello spazio).
    \item \emph{Curva di livello} $c$: $\{(x, y) \in D : f(x, y) = c\}$, cioè l'insieme dei punti del dominio in cui la funzione vale $c$.
\end{itemize}

\begin{esempio}
\leavevmode
\begin{itemize}
    \item \textbf{piano}: $f : \mathbb{R}^2 \to \mathbb{R}$, $f(x, y) = x + y$; le curve di livello $x + y = c$ sono rette;
    \item \textbf{paraboloide}: $f : \mathbb{R}^2 \to \mathbb{R}$, $f(x, y) = x^2 + y^2$; le curve di livello $x^2 + y^2 = c$ (con $c > 0$) sono circonferenze di raggio $\sqrt c$;
    \item \textbf{cono}: $f : \mathbb{R}^2 \to \mathbb{R}$, $f(x, y) = \sqrt{x^2 + y^2}$; le curve di livello $\sqrt{x^2 + y^2} = c$ (con $c > 0$) sono circonferenze di raggio $c$.
\end{itemize}
\end{esempio}

\end{document}
